% GENERATED FILE. Edit canonical lessons, then run npm run generate:books.
% canonical-source-hash: fnv1a64:e3cc2849bfa122d2
% canonical-lessons: ML-C562-talam, ML-C562-ragam, ML-C562-samvidhayakan, ML-C562-tirakkatha, ML-C562-noval

\chapter{Rhythm, Raga and the Screen}
\label{ch:ml-rhythm-raga-and-the-screen}
\hlchaptermodality{562}

\begin{chapteropening}
\textbf{By the end of this chapter:} \emph{I can say rhythm, a raga, a director, a screenplay and a novel.}

Complete the last lesson of chapter 562: I can say rhythm, a raga, a director, a screenplay and a novel.
\end{chapteropening}

\section[\texorpdfstring{\ml{താളം} (tāḷaṁ)}{താളം (tāḷaṁ)}]{\ml{താളം} (tāḷaṁ) — rhythm; a beat}

\label{lesson:ML-C562-talam}



\begin{quote}
Before the new one: say the Malayalam for a guitar, then the Malayalam for a flute.
\end{quote}

\subsection*{You'll want to know: \ml{താളം}}
\textbf{\ml{താളം}} — \emph{tāḷaṁ} — \textquotedblleft{}rhythm; a beat\textquotedblright{}.

\begin{quote}
\ml{കുട്ടികൾ} \ml{പാട്ടിന്റെ} \ml{താളത്തിൽ} \ml{കൈ} \ml{കൊട്ടി}.
\end{quote}

\begin{quote}
\emph{kuṭṭikaḷ pāṭṭinṟe tāḷattil kai koṭṭi.}
\end{quote}

\begin{quote}
The children clapped to the rhythm of the song.
\end{quote}

\ml{താളം} \ml{തെറ്റുക}, to lose the beat, also means to go wrong.

\begin{cousinweb}
\textbf{\ml{താളം}} (\emph{tāḷaṁ}) is borrowed from Sanskrit \emph{tāla}, which CDIAL gives as a cymbal (CDIAL 5801); in Indian music it came to mean the beat the cymbals keep.
\end{cousinweb}

\begin{grammarlens}[title={What you've built}]
One more word for describing what you lived and felt.
\end{grammarlens}

\subsection*{Your turn}
\begin{itemize}
\raggedright
  \item \emph{Say it:} \emph{tāḷaṁ}
  \item \emph{Say it:} \emph{tāḷaṁ}, once more
  \item \emph{Say it:} \emph{kuṭṭikaḷ pāṭṭinṟe tāḷattil kai koṭṭi.} — The children clapped to the rhythm of the song.
  \item \emph{Recall:} say the Malayalam for a guitar, then the Malayalam for a flute, then say \emph{tāḷaṁ} again
\end{itemize}

\begin{tcolorbox}[breakable,colback=teal!4,colframe=teal!35!black,title={Before you move on}]
What does \ml{താളം} mean? (\textquotedblleft{}Rhythm; a beat\textquotedblright{}.) Say it once more.
\end{tcolorbox}
\section[\texorpdfstring{\ml{രാഗം} (rāgaṁ)}{രാഗം (rāgaṁ)}]{\ml{രാഗം} (rāgaṁ) — a raga; a melody}

\label{lesson:ML-C562-ragam}



\begin{quote}
Before the new one: say the Malayalam for a flute, then the Malayalam for rhythm.
\end{quote}

\subsection*{You'll want to know: \ml{രാഗം}}
\textbf{\ml{രാഗം}} — \emph{rāgaṁ} — \textquotedblleft{}a raga; a melody\textquotedblright{}.

\begin{quote}
\ml{ആ} \ml{പാട്ട്} \ml{ഏത്} \ml{രാഗത്തിലാണ്}?
\end{quote}

\begin{quote}
\emph{ā pāṭṭŭ ētŭ rāgattilāṇŭ?}
\end{quote}

\begin{quote}
Which raga is that song in?
\end{quote}

A raga is a set of notes and turns that a melody moves within.

\begin{cousinweb}
\textbf{\ml{രാഗം}} (\emph{rāgaṁ}) is borrowed from Sanskrit \emph{rāga}, \textquotedblleft{}colour, redness\textquotedblright{} (CDIAL 10673); Indian music traditionally explains a raga as what colours the listener's mood.
\end{cousinweb}

\begin{grammarlens}[title={What you've built}]
One more word for describing what you lived and felt.
\end{grammarlens}

\subsection*{Your turn}
\begin{itemize}
\raggedright
  \item \emph{Say it:} \emph{rāgaṁ}
  \item \emph{Say it:} \emph{rāgaṁ}, once more
  \item \emph{Say it:} \emph{ā pāṭṭŭ ētŭ rāgattilāṇŭ?} — Which raga is that song in?
  \item \emph{Recall:} say the Malayalam for a flute, then the Malayalam for rhythm, then say \emph{rāgaṁ} again
\end{itemize}

\begin{tcolorbox}[breakable,colback=teal!4,colframe=teal!35!black,title={Before you move on}]
What does \ml{രാഗം} mean? (\textquotedblleft{}A raga; a melody\textquotedblright{}.) Say it once more.
\end{tcolorbox}
\section[\texorpdfstring{\ml{സംവിധായകൻ} (saṁvidhāyakan)}{സംവിധായകൻ (saṁvidhāyakan)}]{\ml{സംവിധായകൻ} (saṁvidhāyakan) — a director}

\label{lesson:ML-C562-samvidhayakan}



\begin{quote}
Before the new one: say the Malayalam for rhythm, then the Malayalam for a raga.
\end{quote}

\subsection*{You'll want to know: \ml{സംവിധായകൻ}}
\textbf{\ml{സംവിധായകൻ}} — \emph{saṁvidhāyakan} — \textquotedblleft{}a director\textquotedblright{}.

\begin{quote}
\ml{ആ} \ml{സിനിമയുടെ} \ml{സംവിധായകന്} \ml{സമ്മാനം} \ml{കിട്ടി}.
\end{quote}

\begin{quote}
\emph{ā sinimayuṭe saṁvidhāyakanŭ sammānaṁ kiṭṭi.}
\end{quote}

\begin{quote}
The director of that film won a prize.
\end{quote}

\ml{സംവിധാനം} \ml{ചെയ്യുക} is to direct a film.

\begin{cousinweb}
\textbf{\ml{സംവിധായകൻ}} (\emph{saṁvidhāyakan}) is built on Sanskrit \emph{saṁvidhāna}, \textquotedblleft{}arrangement, management\textquotedblright{} (Monier-Williams), with the agent ending -\ml{അകൻ}: the one who arranges.
\end{cousinweb}

\begin{grammarlens}[title={What you've built}]
One more word for describing what you lived and felt.
\end{grammarlens}

\subsection*{Your turn}
\begin{itemize}
\raggedright
  \item \emph{Say it:} \emph{saṁvidhāyakan}
  \item \emph{Say it:} \emph{saṁvidhāyakan}, once more
  \item \emph{Say it:} \emph{ā sinimayuṭe saṁvidhāyakanŭ sammānaṁ kiṭṭi.} — The director of that film won a prize.
  \item \emph{Recall:} say the Malayalam for rhythm, then the Malayalam for a raga, then say \emph{saṁvidhāyakan} again
\end{itemize}

\begin{tcolorbox}[breakable,colback=teal!4,colframe=teal!35!black,title={Before you move on}]
What does \ml{സംവിധായകൻ} mean? (\textquotedblleft{}A director\textquotedblright{}.) Say it once more.
\end{tcolorbox}
\section[\texorpdfstring{\ml{തിരക്കഥ} (tirakkatha)}{തിരക്കഥ (tirakkatha)}]{\ml{തിരക്കഥ} (tirakkatha) — a screenplay; a script}

\label{lesson:ML-C562-tirakkatha}



\begin{quote}
Before the new one: say the Malayalam for a raga, then the Malayalam for a director.
\end{quote}

\subsection*{You'll want to know: \ml{തിരക്കഥ}}
\textbf{\ml{തിരക്കഥ}} — \emph{tirakkatha} — \textquotedblleft{}a screenplay; a script\textquotedblright{}.

\begin{quote}
\ml{അവൾ} \ml{ഒരു} \ml{സിനിമയ്ക്ക്} \ml{തിരക്കഥ} \ml{എഴുതി}.
\end{quote}

\begin{quote}
\emph{avaḷ oru sinimaykkŭ tirakkatha eḻuti.}
\end{quote}

\begin{quote}
She wrote a screenplay for a film.
\end{quote}

\ml{തിര} a screen, a curtain + \ml{കഥ} a story.

\begin{cousinweb}
\textbf{\ml{തിരക്കഥ}} (\emph{tirakkatha}) is a Malayalam compound built from \ml{തിര}, \textquotedblleft{}a curtain\textquotedblright{}, which is native (DEDR 3244), and \ml{കഥ}, from Sanskrit \emph{kathā}, \textquotedblleft{}a story\textquotedblright{} (CDIAL 2704).
\end{cousinweb}

\begin{grammarlens}[title={What you've built}]
One more word for describing what you lived and felt.
\end{grammarlens}

\subsection*{Your turn}
\begin{itemize}
\raggedright
  \item \emph{Say it:} \emph{tirakkatha}
  \item \emph{Say it:} \emph{tirakkatha}, once more
  \item \emph{Say it:} \emph{avaḷ oru sinimaykkŭ tirakkatha eḻuti.} — She wrote a screenplay for a film.
  \item \emph{Recall:} say the Malayalam for a raga, then the Malayalam for a director, then say \emph{tirakkatha} again
\end{itemize}

\begin{tcolorbox}[breakable,colback=teal!4,colframe=teal!35!black,title={Before you move on}]
What does \ml{തിരക്കഥ} mean? (\textquotedblleft{}A screenplay; a script\textquotedblright{}.) Say it once more.
\end{tcolorbox}
\section[\texorpdfstring{\ml{നോവൽ} (nōval)}{നോവൽ (nōval)}]{\ml{നോവൽ} (nōval) — a novel}

\label{lesson:ML-C562-noval}



\begin{quote}
Before the new one: say the Malayalam for a director, then the Malayalam for a screenplay.
\end{quote}

\subsection*{You'll want to know: \ml{നോവൽ}}
\textbf{\ml{നോവൽ}} — \emph{nōval} — \textquotedblleft{}a novel\textquotedblright{}.

\begin{quote}
\ml{ആ} \ml{നോവൽ} \ml{ഞാൻ} \ml{രണ്ട്} \ml{ദിവസം} \ml{കൊണ്ട്} \ml{വായിച്ചുതീർത്തു}.
\end{quote}

\begin{quote}
\emph{ā nōval ñān raṇṭŭ divasaṁ koṇṭŭ vāyiccutīrttu.}
\end{quote}

\begin{quote}
I finished reading that novel in two days.
\end{quote}

One who writes them is a \ml{നോവലിസ്റ്റ്}.

\begin{cousinweb}
\textbf{\ml{നോവൽ}} (\emph{nōval}) was borrowed from English \emph{novel}.
\end{cousinweb}

\begin{grammarlens}[title={What you've built}]
One more word for describing what you lived and felt.
\end{grammarlens}

\subsection*{Your turn}
\begin{itemize}
\raggedright
  \item \emph{Say it:} \emph{nōval}
  \item \emph{Say it:} \emph{nōval}, once more
  \item \emph{Say it:} \emph{ā nōval ñān raṇṭŭ divasaṁ koṇṭŭ vāyiccutīrttu.} — I finished reading that novel in two days.
  \item \emph{Recall:} say the Malayalam for a director, then the Malayalam for a screenplay, then say \emph{nōval} again
\end{itemize}

\begin{tcolorbox}[breakable,colback=teal!4,colframe=teal!35!black,title={Before you move on}]
What does \ml{നോവൽ} mean? (\textquotedblleft{}A novel\textquotedblright{}.) Say it once more.
\end{tcolorbox}
